\documentclass[a4paper]{article} %

\date{} %

% Please, do not change any of the above lines

\usepackage{amssymb} % add other LaTeX packages you need here.
\newtheorem{theorem}{Theorem}
\newtheorem{corollary}[theorem]{Corollary}
\newtheorem{lemma}[theorem]{Lemma}
\newtheorem{proposition}[theorem]{Proposition}
\newtheorem{conjecture}[theorem]{Conjecture}

\begin{document}
\setcounter{page}{1} % Please, do not delete this line

% Your title
\title{Symmetric lattice identities and medians}

% Authors
\author{Author 1, Author 2, Author 3 \\ %
       University of XXYY\\ \\ % Affiliation 1
       Author 4 \\ % If any other author with different Affilation
       University of ZZWW\\ \\ % Affiliation 2 (if needed)
       % New author \\
       % New affiliation \\
       % Add authors and affiliation as needed
       \texttt{author@university.xyz} % Only one corresponding e-mail, please
       }%

\author{Aburub Leen, Gyenizse Gerg\H o.\\
University of Szeged, Szeged, Hungary}

\maketitle

% The abstract

 An n-ary lattice term $f$ is symmetric when its value does not change when we permute its variables i.e, $f(x_1,  \dots, x_n)=f(x_{\sigma(1)},\dots x_{\sigma(n)})$ for all $\sigma \in S_n$. A symmetric lattice identity is an identity where both sides are symmetric terms. An example of a symmetric identity 
 \[
 (x \wedge y)\vee (x \wedge z) \vee (y \wedge z) =(x \vee y)\wedge (x \vee z) \wedge (y \vee z),
 \]
which is famously equivalent to the distributive law (which itself is a nonsymmetric identity).
\begin{conjecture}
    Each lattice variety can be axiomatized by symmetric identities.
\end{conjecture}

In this talk we establish that a stronger version of this conjecture is not true: there is a ternary lattice identity which is not equivalent to any system of ternary symmetric identities.
\begin{thebibliography}{9}
\bibitem{SecondAuthor_BH}
\textsc{G. Gyenizse}, On the symmetric parts of finitely generated free lattices. Acta Mathematica Hungarica, 169(2), 420-431.
\bibitem{Medians}
\textsc{L. Aburub, G. Gyenizse}, Outer and inner medians in some small lattices. arXiv preprint arXiv:2603.17958, (2026).

\end{thebibliography}


\end{document}
